% --- Archon physics formalization source begin ---
% archon:physics
% archon:covers PhyXMiniProblems/problem_phyx_mini_0092.lean
% archon:source-report reports/phyx_mini/problem_phyx_mini_0092.source.json

\chapter{Physics problem phyx\_mini\_0092}
\label{ch:PhyXMiniProblems_problem_phyx_mini_0092}

\paragraph{Problem source.}
\#\# Physical scenario

A laser beam shines along the surface of a block of transparent material (see figure). Half of the beam goes straight to a detector, while the other half travels through the block and then hits the detector. The time delay between the arrival of the two light beams at the detector is 6.25 ns.

\#\# Question

What is the index of refraction of this material?

\#\# Answer choices

- A: A: 1.75

- B: B: 2.12

- C: C: 1.83

- D: D: 1.96

\#\# Figure caption (auxiliary; use the image as primary evidence)

The image illustrates a diagram involving light and refraction. Here's a breakdown:

- **Light Rays:** Three parallel purple lines with arrows depict incident light rays approaching from the left side.
  
- **Medium:** The rays enter a rectangular block labeled with "n = ?" suggesting it is a material of unknown refractive index. The light rays bend as they pass through it.
  
- **Dimensions:** The length of the block is marked as 2.50 m.

- **Detector:** To the right of the medium, there is a vertical line labeled "Detector" where the refracted rays emerge from the block and are directed.

- **Lines:** The rays are parallel upon entering and exit the medium, suggesting refraction effects.

This setup likely represents a refractive index experiment using light passing through the medium to a detector.

\#\# Dataset metadata

Geometrical Optics; Physical Model Grounding Reasoning, Multi-Formula Reasoning

\paragraph{Recorded answer/context.}
A: A: 1.75

\paragraph{Figure/image path.}
/root/proposal\_for\_physic/science-mango/phyx\_mini\_run/phyx\_data/test\_image/92.png

\paragraph{Formalization target.}
create a compiling Lean file with sorry bodies at `PhyXMiniProblems/problem\_phyx\_mini\_0092.lean`. The Lean declarations must preserve the physical quantities, dimensions or dimensional roles, figure labels, governing-law hypotheses, and final relation expressed by this problem.
Use Mathlib/Physlib names found through LeanExplore where available. If a physics API is missing, introduce faithful local abstractions rather than scalar placeholder aliases.

\begin{theorem}[Physics formalization target]
\label{thm:physics:phyx_mini_0092:target}
\lean{PhyXMiniProblems.ProblemPhyXMini0092.refractiveIndex_is_answer_A}
\uses{def:physics:phyx-mini-0092:phyxminiproblems-problemphyxmini0092-laserblockdelaysetup, def:physics:phyx-mini-0092:phyxminiproblems-problemphyxmini0092-lengthinmeters, def:physics:phyx-mini-0092:phyxminiproblems-problemphyxmini0092-durationinseconds, def:physics:phyx-mini-0092:phyxminiproblems-problemphyxmini0092-speedinmeterspersecond, def:physics:phyx-mini-0092:phyxminiproblems-problemphyxmini0092-matchesproblemandfigure, def:physics:phyx-mini-0092:phyxminiproblems-problemphyxmini0092-hasphysicalparameters, def:physics:phyx-mini-0092:phyxminiproblems-problemphyxmini0092-satisfiesopticaldelaylaws, def:physics:phyx-mini-0092:phyxminiproblems-problemphyxmini0092-matchesanswerchoice}
The assigned autoformalize agent should translate this physics problem into Lean declarations in the covered file, with theorem and lemma proof bodies written as `by sorry`.
\end{theorem}
\begin{proof}
This is an autoformalization task, not a proof task. Produce statements that can later be proved by the physics prover without weakening the physical model.
\end{proof}

% --- Archon declaration topology begin ---
\section{Declaration topology}
\begin{definition}[Optical Length]
  \label{def:physics:phyx-mini-0092:phyxminiproblems-problemphyxmini0092-opticallength}
  \lean{PhyXMiniProblems.ProblemPhyXMini0092.OpticalLength}
  A real-valued physical length, independent of the chosen unit readout
\end{definition}
\begin{proof}
  This is the declared model definition of Optical Length.
\end{proof}
\begin{definition}[Optical Duration]
  \label{def:physics:phyx-mini-0092:phyxminiproblems-problemphyxmini0092-opticalduration}
  \lean{PhyXMiniProblems.ProblemPhyXMini0092.OpticalDuration}
  A real-valued physical duration, independent of the chosen unit readout
\end{definition}
\begin{proof}
  This is the declared model definition of Optical Duration.
\end{proof}
\begin{definition}[Optical Speed]
  \label{def:physics:phyx-mini-0092:phyxminiproblems-problemphyxmini0092-opticalspeed}
  \lean{PhyXMiniProblems.ProblemPhyXMini0092.OpticalSpeed}
  A real-valued physical speed, independent of the chosen unit readout
\end{definition}
\begin{proof}
  This is the declared model definition of Optical Speed.
\end{proof}
\begin{definition}[nanosecond Unit Choices]
  \label{def:physics:phyx-mini-0092:phyxminiproblems-problemphyxmini0092-nanosecondunitchoices}
  \lean{PhyXMiniProblems.ProblemPhyXMini0092.nanosecondUnitChoices}
  Unit choices in which duration readouts are measured in nanoseconds
\end{definition}
\begin{proof}
  This is the declared model definition of nanosecond Unit Choices.
\end{proof}
\begin{definition}[Beam Branch]
  \label{def:physics:phyx-mini-0092:phyxminiproblems-problemphyxmini0092-beambranch}
  \lean{PhyXMiniProblems.ProblemPhyXMini0092.BeamBranch}
  The two parallel branches drawn in the source figure
\end{definition}
\begin{proof}
  This is the declared model definition of Beam Branch.
\end{proof}
\begin{definition}[Figure Landmark]
  \label{def:physics:phyx-mini-0092:phyxminiproblems-problemphyxmini0092-figurelandmark}
  \lean{PhyXMiniProblems.ProblemPhyXMini0092.FigureLandmark}
  Named locations appearing in the source figure and its geometry
\end{definition}
\begin{proof}
  This is the declared model definition of Figure Landmark.
\end{proof}
\begin{definition}[Laser Block Delay Setup]
  \label{def:physics:phyx-mini-0092:phyxminiproblems-problemphyxmini0092-laserblockdelaysetup}
  \lean{PhyXMiniProblems.ProblemPhyXMini0092.LaserBlockDelaySetup}
  \uses{def:physics:phyx-mini-0092:phyxminiproblems-problemphyxmini0092-opticallength, def:physics:phyx-mini-0092:phyxminiproblems-problemphyxmini0092-opticalduration, def:physics:phyx-mini-0092:phyxminiproblems-problemphyxmini0092-opticalspeed, def:physics:phyx-mini-0092:phyxminiproblems-problemphyxmini0092-beambranch, def:physics:phyx-mini-0092:phyxminiproblems-problemphyxmini0092-figurelandmark}
  The physical quantities associated with the split-beam timing experiment. The functions indexed by BeamBranch keep the direct and in-material paths distinct even though the figure gives them the same geometric length and the same destination
\end{definition}
\begin{proof}
  This is the declared model definition of Laser Block Delay Setup.
\end{proof}
\begin{definition}[length In Meters]
  \label{def:physics:phyx-mini-0092:phyxminiproblems-problemphyxmini0092-lengthinmeters}
  \lean{PhyXMiniProblems.ProblemPhyXMini0092.lengthInMeters}
  \uses{def:physics:phyx-mini-0092:phyxminiproblems-problemphyxmini0092-opticallength}
  Read a physical length as a real number of metres
\end{definition}
\begin{proof}
  This is the declared model definition of length In Meters.
\end{proof}
\begin{definition}[duration In Seconds]
  \label{def:physics:phyx-mini-0092:phyxminiproblems-problemphyxmini0092-durationinseconds}
  \lean{PhyXMiniProblems.ProblemPhyXMini0092.durationInSeconds}
  \uses{def:physics:phyx-mini-0092:phyxminiproblems-problemphyxmini0092-opticalduration}
  Read a physical duration as a real number of seconds
\end{definition}
\begin{proof}
  This is the declared model definition of duration In Seconds.
\end{proof}
\begin{definition}[duration In Nanoseconds]
  \label{def:physics:phyx-mini-0092:phyxminiproblems-problemphyxmini0092-durationinnanoseconds}
  \lean{PhyXMiniProblems.ProblemPhyXMini0092.durationInNanoseconds}
  \uses{def:physics:phyx-mini-0092:phyxminiproblems-problemphyxmini0092-opticalduration, def:physics:phyx-mini-0092:phyxminiproblems-problemphyxmini0092-nanosecondunitchoices}
  Read a physical duration as a real number of nanoseconds
\end{definition}
\begin{proof}
  This is the declared model definition of duration In Nanoseconds.
\end{proof}
\begin{definition}[speed In Meters Per Second]
  \label{def:physics:phyx-mini-0092:phyxminiproblems-problemphyxmini0092-speedinmeterspersecond}
  \lean{PhyXMiniProblems.ProblemPhyXMini0092.speedInMetersPerSecond}
  \uses{def:physics:phyx-mini-0092:phyxminiproblems-problemphyxmini0092-opticalspeed}
  Read a physical speed as a real number of metres per second
\end{definition}
\begin{proof}
  This is the declared model definition of speed In Meters Per Second.
\end{proof}
\begin{definition}[Matches Problem And Figure]
  \label{def:physics:phyx-mini-0092:phyxminiproblems-problemphyxmini0092-matchesproblemandfigure}
  \lean{PhyXMiniProblems.ProblemPhyXMini0092.MatchesProblemAndFigure}
  \uses{def:physics:phyx-mini-0092:phyxminiproblems-problemphyxmini0092-laserblockdelaysetup, def:physics:phyx-mini-0092:phyxminiproblems-problemphyxmini0092-lengthinmeters, def:physics:phyx-mini-0092:phyxminiproblems-problemphyxmini0092-durationinnanoseconds}
  Problem-text and primary-figure readouts. Both branches span the block's 2.50 m horizontal extent, carry half of the beam, and end at the detector; the in-material branch arrives 6.25 ns later. No refractive-index value is included here
\end{definition}
\begin{proof}
  This is the declared model definition of Matches Problem And Figure.
\end{proof}
\begin{definition}[Has Physical Parameters]
  \label{def:physics:phyx-mini-0092:phyxminiproblems-problemphyxmini0092-hasphysicalparameters}
  \lean{PhyXMiniProblems.ProblemPhyXMini0092.HasPhysicalParameters}
  \uses{def:physics:phyx-mini-0092:phyxminiproblems-problemphyxmini0092-laserblockdelaysetup, def:physics:phyx-mini-0092:phyxminiproblems-problemphyxmini0092-lengthinmeters, def:physics:phyx-mini-0092:phyxminiproblems-problemphyxmini0092-durationinseconds, def:physics:phyx-mini-0092:phyxminiproblems-problemphyxmini0092-speedinmeterspersecond}
  Positivity conditions for the physical quantities in the experiment
\end{definition}
\begin{proof}
  This is the declared model definition of Has Physical Parameters.
\end{proof}
\begin{definition}[Satisfies Optical Delay Laws]
  \label{def:physics:phyx-mini-0092:phyxminiproblems-problemphyxmini0092-satisfiesopticaldelaylaws}
  \lean{PhyXMiniProblems.ProblemPhyXMini0092.SatisfiesOpticalDelayLaws}
  \uses{def:physics:phyx-mini-0092:phyxminiproblems-problemphyxmini0092-laserblockdelaysetup, def:physics:phyx-mini-0092:phyxminiproblems-problemphyxmini0092-lengthinmeters, def:physics:phyx-mini-0092:phyxminiproblems-problemphyxmini0092-durationinseconds, def:physics:phyx-mini-0092:phyxminiproblems-problemphyxmini0092-speedinmeterspersecond}
  Governing laws for the idealized two-branch timing model. The direct branch propagates at Physlib's vacuum speed of light. In the block, v n = c. Constant-speed kinematics gives t v = L on each branch, and the last field identifies the measured delay with the difference in arrival times. None of these laws assumes the requested numerical value of n
\end{definition}
\begin{proof}
  This is the declared model definition of Satisfies Optical Delay Laws.
\end{proof}
\begin{definition}[Answer Choice]
  \label{def:physics:phyx-mini-0092:phyxminiproblems-problemphyxmini0092-answerchoice}
  \lean{PhyXMiniProblems.ProblemPhyXMini0092.AnswerChoice}
  Labels of the four displayed multiple-choice answers
\end{definition}
\begin{proof}
  This is the declared model definition of Answer Choice.
\end{proof}
\begin{definition}[refractive Index Readout]
  \label{def:physics:phyx-mini-0092:phyxminiproblems-problemphyxmini0092-answerchoice-refractiveindexreadout}
  \lean{PhyXMiniProblems.ProblemPhyXMini0092.AnswerChoice.refractiveIndexReadout}
  \uses{def:physics:phyx-mini-0092:phyxminiproblems-problemphyxmini0092-answerchoice}
  Dimensionless refractive-index readout printed beside each answer label
\end{definition}
\begin{proof}
  This is the declared model definition of refractive Index Readout.
\end{proof}
\begin{definition}[Matches Answer Choice]
  \label{def:physics:phyx-mini-0092:phyxminiproblems-problemphyxmini0092-matchesanswerchoice}
  \lean{PhyXMiniProblems.ProblemPhyXMini0092.MatchesAnswerChoice}
  \uses{def:physics:phyx-mini-0092:phyxminiproblems-problemphyxmini0092-answerchoice}
  Agreement with an index displayed to the nearest hundredth
\end{definition}
\begin{proof}
  This is the declared model definition of Matches Answer Choice.
\end{proof}
\begin{lemma}[refractive Index eq one add delay ratio]
  \label{lem:physics:phyx-mini-0092:phyxminiproblems-problemphyxmini0092-refractiveindex-eq-one-add-delay-ratio}
  \lean{PhyXMiniProblems.ProblemPhyXMini0092.refractiveIndex_eq_one_add_delay_ratio}
  \uses{def:physics:phyx-mini-0092:phyxminiproblems-problemphyxmini0092-laserblockdelaysetup, def:physics:phyx-mini-0092:phyxminiproblems-problemphyxmini0092-lengthinmeters, def:physics:phyx-mini-0092:phyxminiproblems-problemphyxmini0092-durationinseconds, def:physics:phyx-mini-0092:phyxminiproblems-problemphyxmini0092-speedinmeterspersecond, def:physics:phyx-mini-0092:phyxminiproblems-problemphyxmini0092-matchesproblemandfigure, def:physics:phyx-mini-0092:phyxminiproblems-problemphyxmini0092-hasphysicalparameters, def:physics:phyx-mini-0092:phyxminiproblems-problemphyxmini0092-satisfiesopticaldelaylaws}
  The equal-path timing laws determine the material index by n = 1 + c Δt / L, using SI scalar readouts of the dimensionful quantities
\end{lemma}
\begin{proof}
  The conclusion follows from the governing relations and figure readouts named in the statement of refractive Index eq one add delay ratio.
\end{proof}
% --- Archon declaration topology end ---
% --- Archon physics formalization source end ---
